% ============================================================
%  APPENDIX C — REPOSITORY ARCHITECTURE AND CONCEPT MAP
% ============================================================

\subsection{Repository Architecture Overview}
\label{subsec:repo_overview}
% ------------------------------------------------------------

The complete implementation accompanying this paper --- the environment,
the baseline solvers, the verification suite, and the literature-review
pipeline described in Section~\ref{sec:literature} --- is organised as a
single version-controlled repository. The repository separates
\emph{documentation} (this manuscript, the thesis, and the literature
corpus) from \emph{implementation} (the environment, ledger, baselines,
and tests), as summarised below.

\begin{lstlisting}[basicstyle=\ttfamily\scriptsize, breaklines=true]
    ppm_env/
    |-- notation.yaml        parameter dictionary (academic <-> code)
    |-- docs/
    |   |-- article/          LaTeX source of this manuscript
    |   |-- lit_archive/      literature screening pipeline and corpus
    |   |-- mscn/             supplementary thesis documentation
    |   |-- pp_presentation/  thesis defense slides (university deliverable)
    |   `-- thesis_farsi/     Persian thesis manuscript (university deliverable)
    `-- src/
        |-- db/               ledger schema, configs, run records
        |-- env/              core CMDP environment implementation
        `-- tests/            verification suite, baselines, gap analysis
    \end{lstlisting}

The \code{mscn/}, \code{pp\_presentation/}, and \code{thesis\_farsi/}
directories satisfy institutional degree-completion requirements and fall
outside the reproducibility scope of this paper; they are listed here only
for completeness. The \code{notation.yaml} file at the repository root
provides an explicit, machine-readable correspondence between every
academic symbol introduced in Section~\ref{sec:problem} and its
counterpart variable name in the source code, and is intended as the
primary cross-reference for readers wishing to connect the formal model to
the implementation.

\medskip

The literature-review pipeline described in
Section~\ref{subsubsec:prisma} resides in \code{docs/lit\_archive/} and
retains, alongside the processing scripts, the full PDF archive
(\code{lit\_archive\_passed/}, \code{lit\_archive\_failed/}) and the
structured JSON extraction records that drove the synthesis table and
gap-analysis figures reported in Section~\ref{sec:literature}.

% - - - - - - - - - - - - - - - - - - - - - - - - - - - - - -
\subsection{Core Contribution Modules}
\label{subsec:repo_contrib}
% - - - - - - - - - - - - - - - - - - - - - - - - - - - - - -

The modules implementing the formal model of Section~\ref{sec:problem}
and the evaluation of Section~\ref{sec:evaluation} are contained entirely
within \code{src/}. The environment dynamics --- state construction,
EVM/S-curve computation, the FIDIC cash-flow cycle, and the termination
gate --- are implemented in \code{src/env/}, with the environment's public
Gymnasium-compatible interface and its internal support routines
maintained as separate modules within that directory. During training or
evaluation, the environment class is instantiated per episode through the
standard Gymnasium \code{reset()}/\code{step()} interface; no additional
driver component is required.

The SQLite audit ledger introduced in Section~\ref{sec:problem} is
schema-defined and persisted under \code{src/db/}, which also holds the
benchmark portfolio configurations used in Section~\ref{sec:evaluation}
and the run-record databases generated by each executed episode.

The verification and evaluation layer resides in \code{src/tests/},
which bundles four functionally distinct components: the Gymnasium
compliance suite, the baseline and optimal solvers, the benchmark
configuration generator together with its archived reference episodes,
and the post-hoc analysis scripts used to produce the figures and tables
of Section~\ref{sec:evaluation}.

\begin{lstlisting}[basicstyle=\ttfamily\scriptsize, breaklines=true]
    src/tests/
    |-- gym_test.py                 Gymnasium API compliance suite
    |-- gym_certificate.png         compliance certificate (Appendix B)
    |-- baseline_solvers.py         greedy heuristics and optimum solver
    |-- reference_episodes_gen/     benchmark configuration generator
    |-- reference_episodes_record/  archived configs and solutions
    `-- analysis/                   gap-analysis and figure scripts
    \end{lstlisting}

The \code{reference\_episodes\_record/} directory retains the configuration
set and recorded solutions used to validate the environment and to produce
the four benchmark configurations evaluated in Section~\ref{sec:evaluation};
it is intended as a working template for practitioners who wish to
instantiate the environment against their own company data.

% - - - - - - - - - - - - - - - - - - - - - - - - - - - - - -
\subsection{Concept-to-Implementation Traceability Matrix}
\label{subsec:repo_matrix}
% - - - - - - - - - - - - - - - - - - - - - - - - - - - - - -

\Cref{tab:repo_map} maps every concept and promised contribution named in
the main text to its corresponding location in the repository, and is
provided to support independent verification of the claims made
throughout this paper.

\renewcommand{\arraystretch}{1.3}
\begin{longtable}{p{4.1cm} p{1.6cm} p{3.9cm} p{4.0cm}}
\caption{Traceability matrix mapping concepts and promised contributions
to their implementation in the repository.}
\label{tab:repo_map} \\
\toprule
\rowcolor{headbg}
\textcolor{headfg}{\textbf{Concept / Contribution}} &
\textcolor{headfg}{\textbf{Ref.}} &
\textcolor{headfg}{\textbf{Repository Location}} &
\textcolor{headfg}{\textbf{Role}} \\
\midrule
\endfirsthead

\multicolumn{4}{c}{{\tablename\ \thetable{} -- continued}} \\
\toprule
\rowcolor{headbg}
\textcolor{headfg}{\textbf{Concept / Contribution}} &
\textcolor{headfg}{\textbf{Ref.}} &
\textcolor{headfg}{\textbf{Repository Location}} &
\textcolor{headfg}{\textbf{Role}} \\
\midrule
\endhead

\midrule
\multicolumn{4}{r}{\textit{continued on next page}} \\
\endfoot

\bottomrule
\endlastfoot

Academic--code notation dictionary & \S\ref{sec:problem} &
\code{notation.yaml} &
Maps formal symbols to code identifiers \\

\rowcolor{RowAlt}
EVM/S-curve state representation &
\S\ref{sec:problem} & \code{src/env/} &
State vector construction \\

FIDIC cash-flow transition dynamics &
\S\ref{sec:problem} & \code{src/env/} &
Advance, retention, interest, and milestone logic \\

\rowcolor{RowAlt}
Two-condition termination gate &
\S\ref{sec:problem} & \code{src/env/} &
Episode termination evaluation \\

SQLite audit ledger &
\S\ref{sec:problem} & \code{src/db/} &
Schema definition and per-run event logging \\

\rowcolor{RowAlt}
Benchmark portfolio configurations &
\S\ref{sec:evaluation} & \code{src/db/} &
Stored configuration instances for evaluation runs \\

Gymnasium compliance verification &
App.~\ref{app:gym_certificate} & \code{src/tests/} &
API conformance tests and certificate generation \\

\rowcolor{RowAlt}
Greedy heuristic baselines (Reach-Plan, Catchup-T, Catchup-Next-T) &
\S\ref{sec:evaluation} & \code{src/tests/} &
Deterministic baseline policies \\

Manually computed full-information optimum &
\S\ref{sec:evaluation} & \code{src/tests/} &
Upper-bound benchmark solver \\

\rowcolor{RowAlt}
Benchmark configuration generation and reference episodes &
\S\ref{sec:evaluation} & \code{src/tests/} &
Configuration matrix, episode archive, and company-data templates \\

Gap-analysis figures and dashboards &
\S\ref{sec:evaluation} & \code{src/tests/} &
Post-hoc analysis and figure generation \\

\rowcolor{RowAlt}
Literature screening, Stages 1--3 &
\S\ref{subsubsec:prisma} & \code{docs/lit\_archive/} &
Readability triage, relevance filtering, structured extraction \\

Synthesis dashboard and literature map &
\S\ref{subsubsec:prisma} & \code{docs/lit\_archive/} &
Gap matrix, domain charts, synthesis table generation \\

\rowcolor{RowAlt}
Manuscript source &
--- & \code{docs/article/} &
LaTeX source, figures, and tables for this paper \\

\end{longtable}

% - - - - - - - - - - - - - - - - - - - - - - - - - - - - - -
\subsection{Reproducibility Notes}
\label{subsec:repo_reproducibility}
% - - - - - - - - - - - - - - - - - - - - - - - - - - - - - -

The mapping in \Cref{tab:repo_map} reflects a specific snapshot of the
repository at the time of writing; subsequent development may reorganise
individual modules without affecting the directory-level structure
described in this appendix. Access to the repository, including the full
commit history, the benchmark configurations, and the Gymnasium
compliance certificate reproduced in Appendix~\ref{app:gym_certificate},
is available from the corresponding author upon request.